\documentclass[10pt,a4paper]{amsart}
\usepackage[margin=19mm]{geometry}
\usepackage{amsmath,amssymb,mathtools}
\usepackage[scr=rsfs,cal=euler]{mathalfa}
\usepackage{xfrac}
\usepackage[unicode,hidelinks,bookmarks=false]{hyperref}
\newcommand{\ie}{i.e.\ }
\newcommand{\cf}{cf.\ }
\newcommand{\resp}{respectively\ }
\newcommand{\Subsection}{\subsection}
\providecommand{\nobreakdash}{\nobreak\hbox{-}}
\setlength{\emergencystretch}{2em}
\multlinegap0pt
\usepackage{enumitem}
\usepackage{aliascnt}
\usepackage[capitalize,noabbrev]{cleveref}
\numberwithin{equation}{section}
\newtheorem{theorem}{Theorem}[section]
\newaliascnt{lemma}{theorem}
\newtheorem{lemma}[lemma]{Lemma}
\aliascntresetthe{lemma}
\newaliascnt{proposition}{theorem}
\newtheorem{proposition}[proposition]{Proposition}
\aliascntresetthe{proposition}
\theoremstyle{remark}
\newaliascnt{remark}{theorem}
\newtheorem{remark}[remark]{Remark}
\aliascntresetthe{remark}
\newaliascnt{assumption}{theorem}
\newtheorem{assumption}[assumption]{Assumption}
\aliascntresetthe{assumption}
\crefname{assumption}{Assumption}{Assumptions}
\newaliascnt{definition}{theorem}
\newtheorem{definition}[definition]{Definition}
\aliascntresetthe{definition}
\newcommand\norm[1]{\left\Vert#1\right\Vert}
\newcommand\sumk{\sum_{k=1}^N}
\newcommand\sumj{\sum_{j=1}^N}
\newcommand\R{\mathbb{R}}
\newcommand\bO{\mathbf{O}}
\newcommand\bw{\mathbf{w}}
\newcommand\bW{\mathbf{W}}
\newcommand\bx{\mathbf{x}}
\newcommand\br{\mathbf{r}}
\newcommand\bX{\mathbf{X}}
\newcommand\by{\mathbf{y}}
\newcommand\tR{\tilde{R}}
\newcommand\bR{\bar{R}}
\newcommand\tA{\tilde{A}}
\newcommand\hA{\hat{A}}
\newcommand\tbw{\tilde{\mathbf{w}}}
\newcommand\bigO{\mathcal{O}}
\newcommand\E{\mathcal{E}}
\newcommand\rT{\mathrm{T}}
\newcommand\bbE{\mathbb{E}}
\newcommand\gp{\mathrm{gp}}
\newcommand\bal{\mathrm{bal}}
\newcommand\param{\mathrm{par}}
\DeclareMathOperator{\diag}{diag}
\setcounter{section}{1}
\title[Approximability estimate of balanced group convolution]{Balanced Group Convolution: An Improved Group Convolution Based on Approximability Estimates: Theorem 3.6}
\author{Youngkyu Lee, Jongho Park, Chang-Ock Lee}
\date{Source: arXiv:2310.12461v1 (CC BY 4.0); DOI 10.1007/s10044-025-01542-6}
\begin{document}
\maketitle
{\sloppy\noindent\emph{Source and licence.} Balanced Group Convolution: An Improved Group Convolution Based on Approximability Estimates. Version arXiv:2310.12461v1 (CC BY 4.0); DOI 10.1007/s10044-025-01542-6. This material is distributed under the Creative Commons Attribution 4.0 International licence, http://creativecommons.org/licenses/by/4.0/, as recorded in the arXiv OAI metadata for this exact version and shown on the abstract page https://arxiv.org/abs/2310.12461v1. Full rights notice: licenses/arxiv-2310.12461-CC-BY-4.0.txt.\par}\smallskip
\noindent\textit{Editorial scope.} Counted claim: original Theorem 3.6 (source label thm:2, source0.tex lines 561-574), the main theoretical result of the article, which the introduction advertises as the estimate of the approximability of the proposed balanced group convolution BGC, delivered with its complete original proof at source lines 576-612 as the single primary proof body. Required context retained uncounted: Section 2 (Group convolution, source lines 264-317) with the matrix form (2.1) of a standard convolution, its block form (2.2) and the group convolution (2.3); Section 3 (Main results) up to the counted statement, that is the definition (3.1) of BGC with the intergroup mean, the definition (3.2) of the approximability, Assumptions 3.1 and 3.2, Lemma 3.3 (Young-type inequality) and Lemma 3.4 (its K/n improvement) with their complete proofs, and Theorem 3.5 with its complete proof (source lines 496-555), the estimate that the counted proof reuses verbatim for the centred inputs; and the concluding Remark 3.7 (source lines 614-620). Every label referenced inside the retained unit is defined inside it: no cross-reference leaves the unit. Printed numbering is reproduced with the author's own declarations (\textbackslash{}newtheorem{theorem}{Theorem}[section], and the assumption, lemma, proposition, remark and definition environments carried by alias counters of the theorem counter through the aliascnt package, so that each environment keeps its own printed name while sharing one number, \textbackslash{}numberwithin{equation}{section}, and an editorial section-counter offset of one, since the paper's Introduction is not part of the unit), so that the delivered PDF prints Section 2, Section 3, equations (2.1)-(2.3), (3.1), (3.2), Assumptions 3.1 and 3.2, Lemmas 3.3 and 3.4, Theorem 3.5, Theorem 3.6 and Remark 3.7 exactly as the published PDF does. Citation handling: every citation command of this article is the natbib author-year \textbackslash{}citep, which the helper's declared reference map binds only to plain \textbackslash{}cite keys; the four \textbackslash{}citep markers that fall inside the retained spans (source lines 301, 315 and 387, keys krizhevsky2012imagenet, lee2022two, zhang2018shufflenet, glorot10a, he2015delving) are therefore removed under a declared bounded\_source\_comparison receipt, the surrounding sentences are kept verbatim, the delivered reference list is empty because the retained unit contains no citation command, and no bibliography entry is transcribed or redistributed. Slips kept literal: the second estimate of Theorem 3.5 writes \textbackslash{}inf\_{\textbackslash{}bW\_{\textbackslash{}gp}} while bounding the balanced-convolution error \textbackslash{}|\textbackslash{}bW\textbackslash{}bx - \textbackslash{}bW\_{\textbackslash{}bal}\textbackslash{}bx\textbackslash{}|\textasciicircum{}2 (source line 506), the proof of Theorem 3.6 repeats \textbackslash{}inf\_{\textbackslash{}bW\_{\textbackslash{}gp}} in the reused display (source line 599), and the BGC paragraph points to Section 3.3 through the printed construction \textbackslash{}cref{Sec:Main}.3 (source line 347); all three are reproduced without change. No mathematics is written, repaired or re-derived; all mathematics is native text with no image dependency.
\section{Group convolution}
\label{Sec:GC}
In this section, we introduce some basic notations for group convolution to be used throughout this paper.
Let $m$ and $n$ be positive integers and $N$ be a common divisor of $m$ and $n$.
We write $m_N = m/N$ and $n_N = n/N$.

We consider convolutional layers that operate on $n$ channels and produces $m$ channels.
Input $\bx$ of a convolutional layer is written as $\bx = (x_{1}, x_{2}, \dots, x_{n})$, where $x_{j}$ is the $j$th channel of $\bx$.
Similarly, output $\by$ is written as $\by = (y_{1}, y_{2}, \dots, y_{m})$, where $y_{i}$ is the $i$th channel of $\by$.
To represent a standard convolutional layer that takes input $\bx$ and produces output $\by$, we use the following matrix expression:
\begin{equation}
\label{eqn:full}
\begin{bmatrix} y_{1} \\ y_{2} \\ \vdots \\ y_{m} \end{bmatrix}
= \begin{bmatrix} W_{11} & W_{12} & \cdots & W_{1n} \\
    W_{21} & W_{22} & \cdots & W_{2n} \\
    \vdots & \vdots & \ddots & \vdots \\
    W_{m1} & W_{m2} & \cdots & W_{mn} \end{bmatrix}
\begin{bmatrix} x_{1} \\ x_{2} \\ \vdots \\ x_{n} \end{bmatrix},
\end{equation}
where $W_{ij}$ is a matrix that represents the generic convolution that maps $x_{j}$ to $y_{i}$.
We suppose that the channels of $\bx$ and $\by$ are evenly partitioned into $N$ groups.
Namely, let $\bx = (\bx^1, \bx^2, \dots, \bx^N)$ and $\by = (\by^1, \dots, \by^N)$, where
\begin{equation*}
    \bx^k = \begin{bmatrix} x_{(k-1)n_N + 1} \\ x_{(k-1)n_N + 2} \\ \vdots \\ x_{k n_N} \end{bmatrix}, \quad
    \by^k = \begin{bmatrix} y_{(k-1)m_N + 1} \\ y_{(k-1)m_N + 2} \\ \vdots \\ y_{k m_N} \end{bmatrix}.
\end{equation*}
Then~\eqref{eqn:full} is rewritten as the following block form:
\begin{equation}
\label{eqn:full_group}
    \begin{bmatrix} \by^1 \\ \by^2 \\ \vdots \\ \by^N \end{bmatrix}
    = \begin{bmatrix} \bW^{11} & \bW^{12} & \cdots & \bW^{1N} \\
        \bW^{21} & \bW^{22} & \cdots & \bW^{2N} \\
        \vdots & \vdots & \ddots & \vdots \\
        \bW^{N1} & \bW^{N2} & \cdots & \bW^{NN} \end{bmatrix}
    \begin{bmatrix} \bx^1 \\ \bx^2 \\ \vdots \\ \bx^N \end{bmatrix}.
\end{equation}

Group convolution, a computationally efficient version of the convolution, is formed by discarding intergroup connection in~\eqref{eqn:full_group}.
That is, we take only the block-diagonal part to obtain the group convolution, such as
\begin{equation}
\label{eqn:group}
    \begin{bmatrix} \by^1 \\ \by^2 \\ \vdots \\ \by^N \end{bmatrix}
    = \begin{bmatrix} \bW^{11} & \bO & \cdots & \bO \\
        \bO & \bW^{22} & \cdots & \bO \\
        \vdots & \vdots & \ddots & \vdots \\
        \bO & \bO & \cdots & \bW^{NN} \end{bmatrix}
    \begin{bmatrix} \bx^1 \\ \bx^2 \\ \vdots \\ \bx^N \end{bmatrix},
\end{equation}
where $\bO$ denotes the zero matrix.
We observe that the number of parameters in the group convolution~\eqref{eqn:group} is $1/N$ of that of the standard convolution~\eqref{eqn:full_group}.
Moreover, the representation matrix in~\eqref{eqn:group} is block-diagonal, allowing each block on the main diagonal to be processed in parallel, making it computationally more efficient than~\eqref{eqn:full_group}. 
However, when the number of parameters in the group convolution is reduced, the performance of a CNN using group convolutional layers may decrease compared to a CNN using standard convolutional layers.

% Section: Main results
\section{Main results}
\label{Sec:Main}
In this section, we present the main results of this paper, including the proposed BGC and estimates of its approximability and computational cost.

% Subsection: Balanced group convolution
\subsection{Balanced group convolution}
A major drawback of the group convolution, which affects performance, is the lack of intergroup communication since the representation matrix given in~\eqref{eqn:group} is block-diagonal.
The motivation behind the proposed BGC is by adding intergroup communication to the group convolution with low computational cost to improve performance.
To achieve this, we utilize a representative vector of small dimension that captures all the information in the input $\bx$. 
Specifically, we use the mean $\bar{\bx}$ of $\{ \bx^k \}$, i.e.,
\begin{equation*}
    \bar{\bx} = \frac{1}{N}\sum_{k=1}^N \bx^k.
\end{equation*}
Then, the proposed BGC is defined by
\begin{equation}
    \label{eqn:balanced}
    \begin{bmatrix} \by^1 \\ \by^2 \\ \vdots \\ \by^N \end{bmatrix}
    = \begin{bmatrix} \bW^{11} & \bO & \cdots & \bO \\
        \bO & \bW^{22} & \cdots & \bO \\
        \vdots & \vdots & \ddots & \vdots \\
        \bO & \bO & \cdots & \bW^{NN} \end{bmatrix}
    \begin{bmatrix} \bx^1 \\ \bx^2 \\ \vdots \\ \bx^N \end{bmatrix}
    + \begin{bmatrix} \bar{\bW}^1 \\ \bar{\bW}^2 \\ \vdots \\ \bar{\bW}^N \end{bmatrix}
    \bar{\bx},
\end{equation}
where each $\bar{\bW}^k$ is a convolution that operates on $n_N$ channels and produces $m_N$ channels.
That is, BGC is an extension of the group convolution that incorporates an additional structure based on the intergroup mean $\bar{\bx}$, which serves as a balancing factor that utilizes the entire information of $\bx$ to intergroup communication.
Specifically, this additional structure allows BGC to better distribute feature representations across groups.

Since the number of parameters in BGC~\eqref{eqn:balanced} is $2/N$ of that of the standard convolution~\eqref{eqn:full_group}, the computational cost of BGC is significantly lower than that of the standard convolution. We will discuss this further in \cref{Sec:Main}.3.

% Subsection: Approximation estimate
\subsection{Approximability estimate}
Thanks to the additional structure in the proposed BGC, we may expect that it behaves more similarly to the standard convolution than the group convolution.
To provide a rigorous assessment of this behavior, we introduce a notion of approximability and prove that BGC has a better approximability to the standard convolution than the group convolution.

Suppose that the standard convolution $\bW = [\bW^{kl}]$ that maps $N$ groups to $N$ groups~(see~\eqref{eqn:full_group}) and the input $\bx = [\bx^k]$ are drawn from a certain probability distribution.
We define the approximability of the group convolution $\bW_{\gp}$ of the form~\eqref{eqn:group} as
\begin{equation}
    \label{eqn:app_err}
    \bbE_{\bW} \left[ \inf_{\bW_{\gp}} \bbE_{\bx} \left[ \| \bW \bx - \bW_{\gp} \bx \|^2 \right] \right],
\end{equation}
where $\| \cdot \|$ denotes the $\ell^2$-norm.
Namely, the approximability~\eqref{eqn:app_err} measures the $\bW$-expectation of the optimal $\bx$-averaged squared $\ell^2$-error between the output $\bW \bx$ of the standard convolution and the output $\bW_{\gp} \bx$ of the group convolution.
The approximability of BGC is defined in the same manner.

To estimate the approximability~\eqref{eqn:app_err}, we need the following assumptions on the distributions of $\bW$ and $\bx$.

% Assumption
\begin{assumption}
\label{ass:data}
A standard convolution $\bW = [W_{ij}]$ and an input $\bx = [x_{j}]$ are random variables that satisfy the following:
\begin{enumerate}[label=(\roman*)]
    \item $\{ W_{ij} \}$ are identically distributed.
    \item $\{ x_{j} \}$ are independent and identically distributed~(i.i.d.).
    \item $\bW$ and $\bx$ are independent.
\end{enumerate}
\end{assumption}
\cref{ass:data}(i) is essential to handle the off-diagonal parts of the group convolution when estimating~\eqref{eqn:app_err}.
To effectively compensate for the absence of the off-diagonal parts in the group convolution using $\bar{\bx}$, we need \cref{ass:data}(ii).
On the other hand, \cref{ass:data}(iii) is a quite natural assumption since it asserts that the convolution and input are independent of each other.
Furthermore, if we additionally assume to the parameters of the standard convolution, we can get better estimates on the approximability of group convolution and BGC.
This assumption is specified in~\cref{ass:input}.
% Assumption
\begin{assumption}
\label{ass:input}
A standard convolution $\bW = [W_{ij}]$ is a random variable such that $\{ W_{ij} \}$ are i.i.d. with $\bbE [ W_{ij} ] = 0$.
\end{assumption}
The independence of $W_{ij}$ is essential to obtain a sharper bound on the expectation of $\| \bW\bx \|^{2}$.
The condition $\bbE [ W_{ij} ] = 0$ usually occurs when the parameters of the convolutional layers are generated by Glorot or He initialization, which are i.i.d. samples from random variable with zero mean.

The following lemma presents a Young-type inequality for the standard convolution, which plays a key role in the proof of the main theorems.

% Lemma: Young's inequality
\begin{lemma}
\label{lem:Young}
Let $\bW$ be a standard convolution that operates on $n$ channels and produces $m$ channels.
For any $n$-channel input $\bx$, we have
\begin{equation*}
\| \bW \bx \| \leq K^{\frac{1}{2}} \| \bW \|_{\param} \| \bx \|,
\end{equation*}
where $K$ is the number of parameters in a convolutional layer that maps a single channel to a single channel and $\| \bW \|_{\param}$ denotes the $\ell^2$-norm of the vector of parameters of the convolution $\bW$.
\end{lemma}
\begin{proof}
    We first prove the case $m = n =1$.
    Let $W$ be a convolution that maps a single channel to a single channel and $x \in \mathbb{R}^D$ be an input.
    Since the output $Wx$ of the convolution is linear with respect to the parameters of $W$, there exists a matrix $X \in \mathbb{R}^{D \times K}$ such that
    \begin{equation*}
        W x = X w,
    \end{equation*}
    where $w \in \mathbb{R}^K$ is the vector of parameters of $W$.
    Note that each entry of $X$ is also an entry of $x$.
    For each entry $x_d$~($1 \leq d \leq D$) of $x$ and each entry $w_k$~($1 \leq k \leq K$) of $w$, their product $w_k x_d$ appears at most one in the formulation of $W x$.
    Hence, the $k$th column $v_k$ of $X$ satisfies
    \begin{equation}
        \label{lem1:Young}
        \| v_k \| \leq \| x \|.
    \end{equation}
    By the triangle inequality, Cauchy--Schwarz inequality, and~\eqref{lem1:Young}, we obtain
    \begin{equation}
        \label{lem2:Young}
        \begin{split}
        \| W x \| &= \| X w \| \leq \sum_{k=1}^K \| v_k w_k \| \\
        &\leq \left( \sum_{k=1}^K \| v_k \|^2 \right)^{\frac{1}{2}} \left( \sum_{k=1}^K w_k^2 \right)^{\frac{1}{2}}
        \leq K^{\frac{1}{2}} \| W \|_{\param} \| x \|,
        \end{split}
    \end{equation}
    which completes the proof of the case $m = n = 1$.

    Next, we consider the general case.
    As in~\eqref{eqn:full}, we write $\bW = [W_{ij}]$ and $\bx = [x_{j}]$.
    By the triangle inequality,~\eqref{lem2:Young}, and the Cauchy--Schwarz inequality, we obtain
    \begin{align*}
        \| \bW \bx \|^2 &= \sum_{i=1}^m \left\| \sum_{j=1}^n W_{ij} x_{j} \right\|^2
        \leq \sum_{i=1}^m \left( \sum_{j=1}^n \| W_{ij} x_{j} \| \right)^2 \\
        &\leq K \sum_{i=1}^m \left( \sum_{j=1}^n \| W_{ij} \|_{\param} \|x_{j} \| \right)^2 \\
        &\leq K \sum_{i=1}^m \left( \sum_{j=1}^n \| W_{ij} \|_{\param}^2 \cdot \sum_{j=1}^n \| x_{j} \|^2 \right) \\
        &=  K \| \bW \|_{\param}^2 \| \bx \|^2,
    \end{align*}
    which is our desired result.
\end{proof}

A direct consequence of \cref{lem:Young} is that, if $\bW$ and $\bx$ are independent, then we have
\begin{equation}
\label{eqn:Young}
    \bbE_{\bW} \bbE_{\bx} \left[ \| \bW \bx \|^2 \right] \leq K \bbE \left[ \| \bW \|_{\param}^2 \right] \bbE \left[ \| \bx \|^2 \right].
\end{equation}
In the following lemma, we show that the above estimate can be improved up to a multiplicative factor $1/n$ if we additionally assume that $W_{ij}$ has zero mean and that some random variables are independent and/or identically distributed.

% Lemma: Improved Young's inequality
\begin{lemma}
    \label{lem:Young_improved}
    Let $\bW = [W_{ij}]$ be a standard convolution that operates on $n$ channels and produces $m$ channels, and let $\bx = [x_j]$ be an $n$-channel input.
    Assume that $\bW$ and $\bx$ satisfy the following:
    \begin{enumerate}[label=\emph{(\roman*)}]
        \item $\{ W_{ij} \}$ are i.i.d. with $\bbE [W_{ij}] = 0$.
        \item $\{ x_j \}$ are identically distributed.
        \item $\bW$ and $\bx$ are independent.
    \end{enumerate}
    Then we have
    \begin{equation*}
        \bbE_{\bW} \bbE_{\bx} \left[ \| \bW \bx \|^2 \right]
        \leq \frac{K}{n} \bbE \left[ \| \bW \|_{\param}^2 \right] \bbE \left[ \| \bx \|^2 \right].
    \end{equation*}
\end{lemma}
\begin{proof}
We first prove the case $m = 1$.
For $i \neq j$, invoking~(i) yields
\begin{equation}
    \label{lem1:Young_improved}
    \bbE_{\bW} \bbE_{\bx} \left[ x_i^{\rT} W_{1i}^{\rT} W_{1j} x_j \right]
    = \bbE_{\bx} \left[ x_i^\rT  \bbE \left[ W_{1i}^{\rT}\right] \bbE \left[ W_{1j} \right] x_j \right] = 0.
\end{equation}
On the other hand, for $1 \leq j \leq n$, by~(iii) and~\eqref{eqn:Young}, we have
\begin{equation}
    \label{lem2:Young_improved}
    \bbE_{W_{1j}} \bbE_{x_j} \left[ \| W_{1j} x_j \|^2 \right]
    \leq K \bbE \left[ \| W_{1j} \|_{\param}^2 \right] \bbE \left[ \|x_j\|^2 \right]
    = \frac{K}{n^2} \bbE \left[ \| \bW \|_{\param}^2 \right] \bbE \left[ \| \bx \|^2 \right],
\end{equation}
where the last equality is due to (i) and (ii).
By~\eqref{lem1:Young_improved} and~\eqref{lem2:Young_improved}, we get
\begin{align*}
    \bbE_{\bW} \bbE_{\bx} \left[ \| \bW \bx \|^2 \right]
    &= \sum_{i=1}^n \sum_{j=1}^n \bbE_{\bW} \bbE_{\bx} \left[ x_i^{\rT} W_{1i}^{\rT} W_{1j} x_j \right] \\
    &\leq \sum_{j=1}^n \frac{K}{n^2} \bbE \left[ \| \bW \|_{\param}^2 \right] \bbE \left[ \| \bx \|^2 \right]
    = \frac{K}{n} \bbE \left[ \| \bW \|_{\param}^2 \right] \bbE \left[ \| \bx \|^2 \right],
\end{align*}
which completes the proof of the case $m = 1$.
The general case $m > 1$ can be shown by applying the case $m = 1$ to each $i$th output channel~($1 \leq i \leq m$) and then summing over all $i$.
\end{proof}

The representation matrix of the group convolution presented in~\eqref{eqn:group} becomes more sparse as the number of groups $N$ increases.
This suggests that the approximability of the group convolution decreases as $N$ increases.
However, to the best of our knowledge, there has been no theoretical analysis on this.
\cref{thm:1} presents the first theoretical result regarding the approximability of the group convolution.

% Theorem: Approximation estimate of the group convolution

\begin{theorem}
  \label{thm:1}
  For an $n$-channel input $\bx$, let $\bW \bx$ and $\bW_{\gp} \bx$ denote the outputs of the standard and group convolutions given in~\eqref{eqn:full_group} and~\eqref{eqn:group}, respectively.
  Under \cref{ass:data}, we have
  \begin{equation*}
  \bbE_{\bW} \left[ \inf_{\bW_{\gp}} \bbE_{\bx} \left[ \| \bW \bx - \bW_{\gp} \bx \|^2 \right] \right]
  \leq K \left( 1 - \frac{1}{N} \right)^2 \bbE \left[ \| \bW \|_{\param}^2 \right] \bbE \left[ \| \bx \|^2 \right].
  \end{equation*}
  In addition, if we further assume that \cref{ass:input} holds, then we have
  \begin{equation*}
  \bbE_{\bW} \left[ \inf_{\bW_{\gp}} \bbE_{\bx} \left[ \| \bW \bx - \bW_{\bal} \bx \|^2 \right] \right]
  \leq \frac{K}{n} \left( 1 - \frac{1}{N} \right) \bbE \left[ \| \bW \|_{\param}^2 \right] \bbE \left[ \| \bx \|^2 \right].
  \end{equation*}
\end{theorem}

% Proof of thm:1
\begin{proof}
Take any standard convolution $\bW = [\bW^{kl}]$ and any input $\bx = [\bx^k]$ as in~\eqref{eqn:full_group}.
If we set
\begin{equation*}
    \bW_{\gp} = \begin{bmatrix} \bW^{11} & \bO & \cdots & \bO \\
        \bO & \bW^{22} & \cdots & \bO \\
        \vdots & \vdots & \ddots & \vdots \\
        \bO & \bO & \cdots & \bW^{NN} \end{bmatrix},
\end{equation*}
i.e., if we choose $\bW_{\gp}$ as the block-diagonal part of $\bW$, then we have
\begin{align}
    \label{thm1:1}
    \bbE_{\bW} \left[ \inf_{\bW_{\gp}} \bbE_{\bx} \left[ \| \bW \bx - \bW_{\gp} \bx \|^2 \right] \right]
    &\leq \bbE_{\bW} \bbE_{\bx} \left[ \| \bW \bx - \bW_{\gp} \bx \|^2 \right] \nonumber \\
    &= \sum_{k=1}^N \bbE_{\bW} \bbE_{\bx} \left[ \left\| \sum_{l \neq k} \bW^{kl} \bx^l \right\|^2 \right].
\end{align}
For each $k$, since the map $\bx \mapsto \sum_{l \neq k} \bW^{kl} \bx^l$ is a convolution that operates on $(n - n/N)$ channels~($\bx_k$ is excluded from the input), invoking \cref{lem:Young,lem:Young_improved} yields
\begin{equation}
    \label{thm2:1}
    \bbE_{\bW} \bbE_{\bx} \left[ \left\| \sum_{l \neq k} \bW^{kl} \bx^l \right\|^2 \right] \leq C_{n,N,k} \left( \sum_{l \neq k} \bbE \left[ \| \bW^{kl} \|_{\param}^2 \right] \right) \left( \sum_{l \neq k} \bbE \left[ \| \bx^l \|^2 \right] \right),
\end{equation}
where
\begin{equation}
\label{C}
C_{n,N,k} = \begin{cases}
K, & \quad \text{ under \cref{ass:data},} \\
\frac{K}{n-n/N}, & \quad \text{ under \cref{ass:data,ass:input}.}
\end{cases}
\end{equation}
Note that the independence condition between $\bW$ and $\bx$ are used in~\eqref{thm2:1}.

Meanwhile, \cref{ass:data} implies that
\begin{equation}
    \label{thm3:1}
    \bbE \left[ \| \bW^{kl} \|_{\param}^2 \right] = \frac{1}{N^2} \bbE \left[ \| \bW \|_{\param}^2 \right], \quad
    \bbE \left[ \| \bx^l \|^2 \right] = \frac{1}{N} \bbE \left[ \| \bx \|^2 \right].
\end{equation}
Finally, combination of~\eqref{thm1:1},~\eqref{thm2:1}, and~\eqref{thm3:1} yields
\begin{equation*}
    \bbE_{\bW} \left[ \inf_{\bW_{\gp}} \bbE_{\bx} \left[ \| \bW \bx - \bW_{\gp} \bx \|^2 \right] \right]
    \leq C_{n,N,k} \left( 1 - \frac{1}{N} \right)^2 \bbE \left[ \| \bW \|_{\param}^2 \right] \bbE \left[ \| \bx \|^2 \right],
\end{equation*}
which completes the proof.
\end{proof}

\begin{theorem}
  \label{thm:2}
  For an $n$-channel input $\bx$, let $\bW \bx$ and $\bW_{\bal} \bx$ denote the outputs of the standard convolution and BGC given in~\eqref{eqn:full_group} and~\eqref{eqn:balanced}, respectively.
  Under \cref{ass:data}, we have
  \begin{equation*}
  \bbE_{\bW} \left[ \inf_{\bW_{\bal}} \bbE_{\bx} \left[ \| \bW \bx - \bW_{\bal} \bx \|^2 \right] \right]
  \leq K \left( 1 - \frac{1}{N} \right)^3 \bbE \left[ \| \bW \|_{\param}^2 \right] \bbE \left[ \| \bx \|^2 \right].
  \end{equation*}
  In addition, if we further assume that \cref{ass:input} holds, we have
  \begin{equation*}
  \bbE_{\bW} \left[ \inf_{\bW_{\bal}} \bbE_{\bx} \left[ \| \bW \bx - \bW_{\bal} \bx \|^2 \right] \right]
  \leq \frac{K}{n} \left( 1 - \frac{1}{N} \right)^{2} \bbE \left[ \| \bW \|_{\param}^2 \right] \bbE \left[ \| \bx \|^2 \right].
  \end{equation*}
\end{theorem}

\begin{proof}
Take any standard convolution $\bW = [\bW^{kl}]$ and any input $\bx = [\bx^k]$ as in~\eqref{eqn:full_group}.
If we choose $\bW_{\bal}$ as
\begin{equation*}
    \bW_{\bal} \bx
    =
    \begin{bmatrix} \bW^{11} & \bO & \cdots & \bO \\
        \bO & \bW^{22} & \cdots & \bO \\
        \vdots & \vdots & \ddots & \vdots \\
        \bO & \bO & \cdots & \bW^{NN} \end{bmatrix}
    \begin{bmatrix} \bx^1 \\ \bx^2 \\ \vdots \\ \bx^N \end{bmatrix}
    + \begin{bmatrix} \sum_{l \neq 1} \bW^{1l} \\ \sum_{l \neq 2} \bW^{2l} \\ \vdots \\ \sum_{l \neq N} \bW^{Nl} \end{bmatrix}
    \bar{\bx},
\end{equation*}
then we get
\begin{equation*}
    \| \bW \bx - \bW_{\bal} \bx \|^2 = \sum_{k=1}^N \left\| \sum_{l \neq k} \bW^{kl} (\bx^l - \bar{\bx}) \right\|^2.
\end{equation*}
Note that, in the proof of~\cref{thm:1}, the independence condition of $\{ x_j \}$ was never used.
Hence, we can use the same argument as in the proof of \cref{thm:1} to obtain
\begin{equation}
\label{thm1:2}
    \begin{split}
        &\bbE_{\bW} \left[ \inf_{\bW_{\gp}} \bbE_{\bx} \left[ \| \bW \bx - \bW_{\gp} \bx \|^2 \right] \right] \\
        &\leq C_{n,N,k} \left( 1 - \frac{1}{N} \right)^2 \bbE \left[ \| \bW \|_{\param}^2 \right] \bbE \left[ \| \bx - \bar{\bx }\|^2 \right] \\
        &= C_{n,N,k} \left( 1 - \frac{1}{N} \right)^2 \bbE \left[ \| \bW \|_{\param}^2 \right] \left( \bbE \left[ \| \bx \|^2 \right] - N \bbE \left[ \| \bar{\bx} \|^2 \right] \right),
    \end{split}
\end{equation}
where $C_{n,N,k}$ was given in~\eqref{C}.
Meanwhile, since $\{ \bx^k \}$ are i.i.d., we have
\begin{equation}
\label{thm2:2}
    \bbE \left[ \| \bar{\bx} \|^2 \right] = \frac{1}{N^2} \bbE \left[ \| \bx \|^2 \right] + \frac{N-1}{N} \left\| \bbE \left[ \bx^1 \right] \right\|^2
    \geq \frac{1}{N^2} \bbE \left[ \| \bx \|^2 \right].
\end{equation}
Combining~\eqref{thm1:2} and~\eqref{thm2:2} yields the desired result.
\end{proof}

\begin{remark}
By \cref{ass:data}, the approximability of the group convolution has a bound of $K(1-1/N)^{2}$.
Adding \cref{ass:input} to this provides $K(1-1/N)/n$ bound.
Since $n \geq N$, the latter is a sharper bound.
On the other hand, in both cases, the approximability of BGC has sharper estimates than the group convolution by $(1-1/N)$ factor.
\end{remark}


\end{document}
